% ch6-storage.tex
\chapter{From Immutable Values to Mutable Storage}
\label{ch:storage}

\begin{answer}
Kernel selection produces BLAS and LAPACK kernels on \emph{tensors}, in
destination-passing style: each kernel names the value it overwrites, and
returns the new value. MLIR's One-Shot Bufferize then assigns buffers. It
lets a kernel overwrite its destination in place when the destination's
buffer may be written and no later use reads the old value; otherwise it
inserts a copy. Read-only parameters are protected by an attribute, so the
caller's \(M\) is copied before \code{potrf} overwrites it, unless the
caller declares it overwritable. After bufferization, ulm.lam's own pass
places all temporaries in one caller-provided workspace at disjoint
offsets. No type system is involved: in-place updates are justified by an
analysis, and the absence of aliasing by construction.
\end{answer}

% ---------------------------------------------------------------------------
\section{Destination-passing style kernels}
\label{sec:storage-dps}

\begin{steps}
\item[Input] a kernel that BLAS executes in place, \(B \leftarrow L\inv B\).
\item[Representation] \code{ulm_lamk.trsm ... ins(\%L) outs(\%B) -> tensor} on tensors; the same op with no result on memrefs.
\item[Rule] the \code{outs} operand is the DPS ``init'': on tensors, result \(i\) is the new value of init \(i\) and has its type (\code{verifyDps}).
\item[Output] a value-semantic op whose natural buffer is that of its \code{outs} operand.
\item[Theory] \FG{10.3}{Bufferization and destination-passing style}.
\end{steps}

All kernels share one ODS base class
(\srcref{include/ulm_lam/UlmLamKernelOps.td}{30}{45}):

\begin{lstlisting}[style=td]
class UlmLamk_Op<string mnemonic, string init, list<Trait> traits = []>
    : Op<UlmLamk_Dialect, mnemonic, !listconcat(traits, [
        DestinationStyleOpInterface,
        DeclareOpInterfaceMethods<MemoryEffectsOpInterface>,
        DeclareOpInterfaceMethods<UlmLam_PropertyInferenceOpInterface>,
        DeclareOpInterfaceMethods<UlmLamk_KernelCostOpInterface>])> {
  let results = (outs Variadic<2DTensorOf<[F32, F64]>>:$result);
  ...
}
\end{lstlisting}

The variadic result is the trick that lets one op serve both worlds: on
tensors it has one result (the new value), on memrefs none. The GLS
\code{trsm}, before and after bufferization:

\begin{lstlisting}
%3 = ulm_lamk.trsm side(left) uplo(lower) trans(none) diag(non_unit) alpha(%cst_0)
       ins(%result : tensor<?x?xf64>) outs(%2 : tensor<?x?xf64>) -> tensor<?x?xf64>
ulm_lamk.trsm side(left) uplo(lower) trans(none) diag(non_unit) alpha(%cst)
       ins(%view : memref<?x?xf64>) outs(%view_2 : memref<?x?xf64>)
\end{lstlisting}

The bufferization model for every kernel
(\srcref{lib/UlmLamDialect.cpp}{1168}{1213}) derives from MLIR's
\code{DstBufferizableOpInterfaceExternalModel}, which supplies the analysis
(which operands are read, which are written, which result aliases which
destination). ulm.lam implements only \code{bufferize()}: create the same
op on buffers and replace each tensor result by its destination's buffer.
The comment says it in one line: ``Analysis (in-place decisions, copies)
comes from the DPS default model.''

% ---------------------------------------------------------------------------
\section{Who may be overwritten}
\label{sec:storage-intent}

One-Shot Bufferize treats function arguments as writable unless told
otherwise. \code{ulm-lam-results-to-out-params} tells it
(\srcref{lib/UlmLamResultsToOutParams.cpp}{344}{368}): a parameter marked
\code{ulm_lam.in} (the default) gets \code{bufferization.writable = false},
one marked \code{ulm_lam.inout} gets \code{true}. Results become out
parameters that the caller allocates.

\begin{center}\small
\begin{tabularx}{\linewidth}{@{}P{3.6cm}P{2.6cm}L@{}}
\toprule
\textbf{Variant} & \textbf{Copies} & \textbf{What changes} \\
\midrule
\code{gls}: \(X\), \(M\), \(y\) \code{in} & 3 (\(M\), \(X\), \(y\)) &
  \code{potrf} and \code{trsv} work on copies in the workspace \\
\code{gls_inout}: \(M\), \(y\) \code{inout} & 1 (\(X\)) & \code{potrf
  uplo(lower) outs(\%arg1 ...)} works on the caller's \(M\); the caller's
  \(M\) and \(y\) are destroyed, as with LAPACK \\
\bottomrule
\end{tabularx}
\end{center}

Both lines are checked by \code{examples/gls/check.sh}
(\cref{app:repro}). In Julia, \code{InOut(Mat(...))} declares the second
kind of parameter. The intent is \Odecl{} by the programmer; that the
generated code respects it is \Oana{} (One-Shot Bufferize's conflict
analysis honors \code{writable = false}).

% ---------------------------------------------------------------------------
\section{Loop state in place}
\label{sec:storage-loops}

A loop-carried matrix should live in one buffer for the whole loop.
One-Shot Bufferize can arrange that only if the value yielded at the end of
the body is computed in the buffer of the iteration argument. Kernel
selection makes sure of it (\srcref{lib/UlmLamPasses.cpp}{358}{377}): if
the yielded value was computed elsewhere, it inserts a
\code{bufferization.materialize_in_destination} into the iteration
argument. In \code{iterate} (\code{outputs/iterate-selected.mlir}):

\begin{lstlisting}
%0 = scf.for %arg4 = %c0 to %arg3 step %c1 iter_args(%arg5 = %arg2) -> (tensor<?x1xf64>) {
  %1 = ulm_lamk.trsv uplo(lower) trans(none) diag(non_unit) ins(%result : ...) outs(%arg5 : ...) -> ...
  %2 = ulm_lamk.trsv uplo(lower) trans(trans) diag(non_unit) ins(%result : ...) outs(%1 : ...) -> ...
  %3 = bufferization.materialize_in_destination %2 in %arg5 : (...) -> tensor<?x1xf64>
  scf.yield %3 : tensor<?x1xf64>
}
\end{lstlisting}

Here both solves already write into \code{\%arg5}'s buffer, and the
materialization costs nothing. For \(x \leftarrow Fx\), which cannot be
computed in place by \code{gemv}, it becomes one copy per iteration instead
of one allocation per iteration.

% ---------------------------------------------------------------------------
\section{One workspace for all temporaries}
\label{sec:storage-workspace}

\begin{steps}
\item[Input] the bufferized function, with a \code{memref.alloc} per temporary.
\item[Representation] after \code{ulm-lam-workspace}: two extra parameters, a byte buffer and a size/protocol word; every temporary a \code{memref.view} into the buffer.
\item[Rule] offsets are the running sum of the temporaries' sizes, each rounded up to 64 bytes; the size is a formula in the run-time sizes.
\item[Output] no allocation inside the function (unless the caller asks for it).
\item[Theory] an arena allocator with a LAPACK-style \code{lwork} protocol; buffer assignment as storage allocation, \FG{10.3}{Bufferization and destination-passing style}, by analogy with register assignment, \FG{4.1}{SSA destruction: parallel copies, lost-copy and swap problems}.
\end{steps}

The protocol (\srcref{lib/UlmLamWorkspace.cpp}{1}{47}): with
\code{*workspace_info == -1} the function only stores the required size and
returns; with \code{0} it allocates the block itself and frees it at the end;
with a positive value it uses the given block, or returns
\(\mathrm{info} = -i\) if the block is too small. The Julia runtime queries
the size once per combination of run-time sizes and keeps one workspace per
task, so repeated calls allocate nothing.

For GLS the bufferized IR computes four aligned sizes and their offsets
(the \code{arith} ops at the top of \texttt{6-bufferized.mlir}).
\Cref{fig:workspace} shows the layout and the numbers for the example of
\cref{sec:gls-numbers}.

\begin{figure}[h]
\centering
\begin{tikzpicture}[font=\sffamily\small,
  seg/.style={draw=accent,fill=accent!8,minimum height=0.9cm,align=center,inner sep=2pt},
  lbl/.style={font=\sffamily\scriptsize,text=metacolor}]
  \node[seg,minimum width=5.2cm,anchor=west] (a) at (0,0) {copy of \(M\), then \(L_1\)\\[-1pt]\scriptsize\(8n^2\)};
  \node[seg,minimum width=2.0cm,anchor=west,fill=oinf!12] (b) at (a.east) {\(S\), then \(L_2\)\\[-1pt]\scriptsize\(8m^2\)};
  \node[seg,minimum width=3.4cm,anchor=west] (c) at (b.east) {\(W\)\\[-1pt]\scriptsize\(8nm\)};
  \node[seg,minimum width=2.2cm,anchor=west] (d) at (c.east) {copy of \(y\), \(z\)\\[-1pt]\scriptsize\(8n\)};
  \node[lbl,anchor=north] at (a.south west) {0};
  \node[lbl,anchor=north] at (b.south west) {\code{\%4}};
  \node[lbl,anchor=north] at (c.south west) {\code{\%9}};
  \node[lbl,anchor=north] at (d.south west) {\code{\%13}};
  \node[lbl,anchor=north] at (d.south east) {\code{\%16}};
  \node[lbl,anchor=north west,align=left] at (0,-1.0) {\(n=4,\ m=2\): offsets 0, 128, 192, 256; total 320 bytes.\\
    \(n=2000,\ m=50\): offsets 0, 32\,000\,000, 32\,020\,032, 32\,820\,032; total 32\,836\,032 bytes.};
\end{tikzpicture}
\caption{The GLS workspace: one block, four views at 64-byte aligned
offsets (each size rounded up to a multiple of 64). \(\beta\) is not in it:
it is written directly into the caller's output.}
\label{fig:workspace}
\end{figure}

Every temporary has its own region of the block for the whole call. The
pass does not overlap temporaries with disjoint lifetimes (in GLS that
would not help: while \(z\) is computed, \(L_1\), \(L_2\), \(W\) and
\(z\) are all live). It does let the branches of an \code{scf.if} share
their part, because only one of them runs. Reuse happens \emph{across}
calls, through the caller's workspace.

% ---------------------------------------------------------------------------
\section{Aliasing}
\label{sec:storage-alias}

BLAS and LAPACK assume that an output does not overlap an input (Fortran's
rules). ulm.lam guarantees this by construction rather than by analysis:
\begin{itemize}
\item temporaries are views at disjoint offsets of one block;
\item read-only parameters are never written (\code{writable = false});
\item a kernel reads its \code{ins} and overwrites only its \code{outs}, and
  bufferization gives an \code{outs} operand that is still needed elsewhere
  a copy.
\end{itemize}

MLIR's generic passes see less. The memref effects of \code{ulm_lamk} ops
are all on the \code{DefaultResource} (\cref{sec:effects}), so any two
kernels on memrefs are treated as possibly touching the same memory. That
is conservative and therefore sound; it costs nothing here, because all
optimization happens before bufferization, on values.

\begin{theory}[ (\FGs{5.5}, \FGs{2.7})]
\FG{5.5}{Alias and points-to analysis: Andersen, Steensgaard, type-based}
answers ``may these two pointers refer to the same memory?''. ulm.lam
avoids the question: on tensors there is no memory and hence no aliasing
(\FGs{2.7}), and after bufferization the program is no longer optimized. The
one analysis that does reason about aliasing is One-Shot Bufferize's, which
tracks which tensor values are equivalent to which buffers in order to
decide in-place writes.
\end{theory}

\begin{limit}
Bufferization and storage reuse do not imply linear typing. A linear type
system (\FGs{1.2}) or a uniqueness type (\FGs{1.4}) would \emph{guarantee},
by a typing rule, that a value is not used after it is updated in place.
Here a value may be used any number of times, and the bufferizer
\emph{decides}, for each update, whether an in-place write is safe; if not,
it copies. The program is never rejected for using a value twice; it only
gets slower.
\end{limit}

% ---------------------------------------------------------------------------
\section{What each lowering keeps and drops}
\label{sec:storage-lost}

\begin{table}[h]
\centering\small
\begin{tabularx}{\linewidth}{@{}P{2.7cm}P{2.6cm}P{2.6cm}P{2.6cm}L@{}}
\toprule
\textbf{Information} & \textbf{ulm\_lam} & \textbf{ulm\_lamk, tensors} & \textbf{ulm\_lamk, memrefs} & \textbf{C++} \\
\midrule
matrix properties & encodings & kernel choice, \code{uplo}, \code{diag}; encodings until \code{strip-properties} & gone & routine names (\code{trlsv}, \code{trusv}) \\
transposition & \code{transpose} ops & \code{trans} flags & \code{trans} flags & swapped strides \\
the algebra & \code{mul}, \code{inverse} & kernels & kernels & calls \\
value semantics & tensors & tensors (DPS) & buffers & pointers \\
run-time sizes & \code{?} + \code{ulm_lam.dims} & \code{index} values & \code{index} values & \code{size_t} arguments \\
in/out intent & arg.\ attributes & arg.\ attributes & copies or none & interface comment \\
failure & --- & \code{info} result & \code{info}, then \code{scf.if} & return value \\
\bottomrule
\end{tabularx}
\caption{Information through the stages. Properties are consumed by kernel
selection: after it, the program no longer says \emph{why} a Cholesky
factorization was chosen, only \emph{that} it was.}
\label{tab:lost}
\end{table}

\Cref{tab:lost} answers ``what is preserved or lost during lowering''. The
pattern is that of \FG{3.3}{Progressive lowering and the dialect conversion
framework}: each stage keeps exactly what a later stage still needs. The
properties are needed only by kernel selection and are dropped right after
it; the DPS ties are needed by bufferization and disappear with it;
transposition survives down to the C++ call, as strides.
